% Section 12.8, from lectures/L12/appendix/b_exercises.md.

\section{Uppgifter}
\label{c12:sec:uppgifter}

\begin{aside}
\textbf{Så arbetar du med uppgifterna.} Del~1 och del~2 räknas under lektionen: först på egen
hand, några minuter per uppgift, sedan gemensamt. Del~3 är extrauppgifter att träna vidare på
efter lektionen; de flesta av dem går att räkna i huvudet eller med papper och penna.

\facitnotis{L12}
\end{aside}

% ------------------------------------------------------------------------------------------
\uppgiftsdel{Del 1}{Repetitionsuppgifter}

\uppgift{1.1}{Linjär förstärkning}
Sambandet mellan den linjära förstärkningen och förstärkningen mätt i dB är
\[
  G_{\text{dB}} = 20\log_{10} G_{\text{lin}},
\]
där $G_{\text{dB}}$ är förstärkningen i dB och $G_{\text{lin}}$ den linjära förstärkningen. En
ljudförstärkare har en förstärkning på $26\,\text{dB}$. Beräkna motsvarande linjära
spänningsförstärkning.

\uppgift{1.2}{Uteffekt i dBm samt mW}
Sambandet mellan en effekt $P$ och motsvarande effekt i dBm (dB i förhållande till
$1\,\text{mW}$) är
\[
  P_{\text{dBm}} = 10\log_{10}\frac{P}{1\,\text{mW}},
\]
där $P$ är effekten i mW och $P_{\text{dBm}}$ motsvarande effekt i dBm. En radiosändare har en
uteffekt på $17\,\text{dBm}$. Ange uteffekten i mW.

\uppgift{1.3}{Ljudnivå från multipla ljudkällor}
För ett godtyckligt antal lika starka ljudkällor kan den totala ljudnivån beräknas med
\[
  L_{\text{tot}} = L_{\text{enkel}} + 10\log_{10} n,
\]
där $L_{\text{tot}}$ är den totala ljudnivån, $L_{\text{enkel}}$ ljudnivån från en ljudkälla och
$n$ antalet ljudkällor. Beräkna antalet lika starka ljudkällor om ljudet från en ljudkälla är
$68\,\text{dB}$ och den totala ljudnivån uppgår till $72{,}8\,\text{dB}$.

% ------------------------------------------------------------------------------------------
\needspace{12\baselineskip}
\uppgiftsdel{Del 2}{Nytt stoff}

\uppgift{2.1}{Derivering av funktioner}
Derivera följande funktioner.
\begin{delar}(2)
  \task $f(x) = -3x^2 + 5x + 2$
  \task $f(x) = 2x^3 - 5x^2 + \dfrac{2x}{3} - 4$
  \task* $f(x) = 2x^4 - x^3 + \dfrac{3x^2}{5} + 4x - 1$
\end{delar}

\uppgift{2.2}{Analys av en parabel}
Betrakta funktionen
\[
  f(x) = -x^2 + 8x - 7.
\]
\begin{parts}
  \item Derivera funktionen $f(x)$, det vill säga bestäm uttrycket för $f'(x)$.
  \item Bestäm var funktionen är stationär, det vill säga lös ekvationen $f'(x) = 0$.
  \item Avgör med hjälp av andra derivatan om punkten är ett \textbf{maximum} eller ett
    \textbf{minimum}.
  \item Beräkna funktionens största eller minsta värde.
  \item Rita grafen till $f(x)$ i GeoGebra, \url{https://www.geogebra.org/graphing?lang=en}, och
    kontrollera att dina svar stämmer.
\end{parts}

\uppgift{2.3}{Analys av strömförlopp i en RC-krets}
Strömmen genom en viss RC-krets kan approximeras med polynomfunktionen
\[
  i(t) = -0{,}2t^3 + 1{,}8t^2 - 3{,}6t + 2{,}5,
\]
där $i(t)$ är strömmen genom kretsen i A och $t$ tiden i sekunder.
\begin{parts}
  \item Derivera funktionen för att bestämma uttrycket för strömmens förändring
    (strömändringshastigheten) $i'(t)$.
  \item Beräkna den tidpunkt (eller de tidpunkter) då strömmen är stationär, det vill säga när
    $i'(t) = 0$.
  \item Bestäm med hjälp av \textbf{andra derivatan} om respektive stationär punkt är ett
    \textbf{maximum} eller ett \textbf{minimum}.
  \item Beräkna strömmens maximi- och minimivärde, det vill säga strömmen $i(t)$ i de
    stationära punkterna.
  \item Rita grafen till $i(t)$ i GeoGebra, \url{https://www.geogebra.org/graphing?lang=en}, och
    kontrollera att dina svar stämmer.
\end{parts}

% ------------------------------------------------------------------------------------------
\uppgiftsdel{Del 3}{Extrauppgifter}
Uppgifterna nedan är extra träning och görs med fördel på egen hand efter lektionen.

\uppgift{3.1}{Derivera polynom}
Derivera följande funktioner.
\begin{delar}(3)
  \task $f(x) = x^5$
  \task $f(x) = 4x^3$
  \task $f(x) = 7$
  \task $f(x) = 2x^2 - 5x + 1$
  \task $f(x) = -3x^4 + x$
  \task $f(x) = \dfrac{x^3}{3}$
\end{delar}

\uppgift{3.2}{Ändringskvot och derivatans definition}
Funktionen $f(x) = x^2$ är given.
\begin{parts}
  \item Ställ upp ändringskvoten över intervallet $[2,\, 2 + h]$ och förenkla den.
  \item Vilket värde närmar sig kvoten när $h \to 0$?
  \item Kontrollera svaret med deriveringsregeln $f'(x) = 2x$.
  \item Beräkna ändringskvoten för $h = 0{,}1$ och $h = 0{,}01$.
\end{parts}

\uppgift{3.3}{Derivatan i en punkt}
Funktionen $f(x) = x^3 - 3x$ är given.
\begin{delar}(3)
  \task Bestäm $f'(x)$.
  \task Beräkna $f'(0)$.
  \task Beräkna $f'(2)$.
  \task* I vilka punkter är tangenten vågrät?
\end{delar}

\uppgift{3.4}{Tangentens ekvation}
Funktionen $f(x) = x^2 - 4x + 5$ är given. Betrakta punkten där $x_0 = 3$.
\begin{delar}(2)
  \task Beräkna $f(3)$.
  \task Beräkna $f'(3)$.
  \task Bestäm tangentens ekvation.
  \task Var skär tangenten $x$-axeln?
\end{delar}

\uppgift{3.5}{Stationär punkt hos en parabel}
Funktionen $f(x) = x^2 - 6x + 8$ är given.
\begin{parts}
  \item Bestäm $f'(x)$.
  \item Lös $f'(x) = 0$.
  \item Avgör med $f''(x)$ om punkten är ett maximum eller ett minimum.
  \item Beräkna funktionens extremvärde.
\end{parts}

\uppgift{3.6}{Stationära punkter hos en tredjegradsfunktion}
Funktionen $f(x) = x^3 - 3x^2 - 9x + 5$ är given.
\begin{parts}
  \item Bestäm $f'(x)$ och faktorisera uttrycket.
  \item Bestäm de stationära punkternas $x$-värden.
  \item Avgör med $f''(x)$ vilken punkt som är maximum respektive minimum.
  \item Beräkna funktionsvärdena i båda punkterna.
\end{parts}

\uppgift{3.7}{Växande och avtagande}
Funktionen $f(x) = x^2 - 4x$ är given.
\begin{delar}(2)
  \task Bestäm $f'(x)$.
  \task För vilka $x$ är funktionen växande?
  \task För vilka $x$ är den avtagande?
  \task Vad händer vid $x = 2$?
\end{delar}

\uppgift{3.8}{Andra derivatan}
Bestäm $f'(x)$ och $f''(x)$.
\begin{delar}(3)
  \task $f(x) = x^4 - 2x^2$
  \task $f(x) = 5x^3 + 2x$
  \task $f(x) = -x^2 + 7x - 1$
  \task* Vad säger tecknet på $f''(x)$ om kurvans form?
\end{delar}

\uppgift{3.9}{Ström genom en kondensator}
Strömmen genom en kondensator ges av $i_C(t) = C\dfrac{du}{dt}$. En kondensator med
$C = 100\,\mu\text{F}$ har spänningen $u(t) = 4t^2 + 2t$ volt.
\begin{parts}
  \item Bestäm $u'(t)$.
  \item Bestäm ett uttryck för $i_C(t)$ i mA.
  \item Beräkna strömmen vid $t = 1\,\text{s}$.
  \item Vid vilken tidpunkt är strömmen $2\,\text{mA}$?
\end{parts}

\uppgift{3.10}{Spänning över en spole}
Spänningen över en spole ges av $u_L(t) = L\dfrac{di}{dt}$. En spole med $L = 0{,}5\,\text{H}$
genomflyts av strömmen $i(t) = 2t^2 - 4t + 3$ ampere.
\begin{parts}
  \item Bestäm $i'(t)$.
  \item Bestäm ett uttryck för $u_L(t)$.
  \item Vid vilken tidpunkt är $u_L = 0$?
  \item Vad händer med strömmen vid denna tidpunkt?
\end{parts}

\uppgift{3.11}{Maximal effekt}
Effekten som en generator levererar till en last varierar med strömmen enligt
$P(I) = 24I - 3I^2$ watt.
\begin{parts}
  \item Bestäm $P'(I)$.
  \item Vid vilken ström är effekten stationär?
  \item Avgör med $P''(I)$ om det är ett maximum eller ett minimum.
  \item Beräkna den maximala effekten.
\end{parts}

\uppgift{3.12}{Förändringshastighet i en signal}
En spänning ges av $u(t) = -2t^2 + 12t$ volt för $0 \leq t \leq 6\,\text{s}$.
\begin{parts}
  \item Bestäm $u'(t)$.
  \item Vid vilken tidpunkt är spänningen störst, och hur stor är den då?
  \item Hur snabbt ändras spänningen vid $t = 1\,\text{s}$?
  \item Vid vilken tidpunkt minskar spänningen med $4\,\text{V/s}$?
\end{parts}

\uppgift{3.13}{Tolkning av derivatans tecken}
För en funktion $f$ gäller att $f'(1) = 3$, $f'(2) = 0$, $f'(3) = -2$ och $f''(2) = -5$.
\begin{parts}
  \item Är $f$ växande eller avtagande vid $x = 1$?
  \item Är $f$ växande eller avtagande vid $x = 3$?
  \item Vilken typ av punkt är $x = 2$?
  \item Beskriv med ord hur grafen ser ut kring $x = 2$.
\end{parts}

\uppgift{3.14}{Hitta felet}
Varje rad innehåller ett vanligt deriveringsfel. Förklara felet och ange det korrekta svaret.
\begin{delar}(2)
  \task $f(x) = x^3 \quad \Rightarrow \quad f'(x) = 3x^2 \cdot x$
  \task $f(x) = 5 \quad \Rightarrow \quad f'(x) = 5$
  \task $f(x) = 4x \quad \Rightarrow \quad f'(x) = 4x$
  \task $f(x) = 3x^2 + 2x \quad \Rightarrow \quad f'(x) = 6x + 2x$
  \task* Om $f'(x_0) = 0$ är $x_0$ alltid en maximipunkt.
  \task* $f''(x)$ fås genom att kvadrera $f'(x)$.
\end{delar}
